\section{Методика исследования}
\subsection{Среда и инструменты}
\begin{tabularx}{\linewidth}{@{}p{44mm}Y@{}}
\hline
Параметр & Фактическое значение \\
\hline
ОС и среда запуска & \Field{Windows, Linux или WSL; версия} \\
Процессор и RAM & \Field{модель и объём} \\
Компилятор & \Field{название и версия} \\
Обычная сборка & \Field{точная команда и флаги} \\
Профилируемая сборка & \Field{команда; для gprof учесть -pg} \\
Диагностика памяти & \Field{средство, версия и команда} \\
Замер времени/RSS & \Field{средство и единицы} \\
Базовая версия & \Field{коммит и корректность} \\
Исправленная версия & \Field{коммит и изменения} \\
\hline
\end{tabularx}

Планируемый набор: gprof для времени функций \cite{gprof},
Valgrind Memcheck для ошибок памяти \cite{memcheck}, при необходимости
Massif для динамики кучи \cite{massif}. Фактический выбор указать выше.
Время обычной программы измерять без Valgrind и санитайзеров;
диагностику проводить отдельными запусками.

\subsection{Нагрузки и корректность}
\begin{tabularx}{\linewidth}{@{}p{41mm}Y@{}}
\hline
Нагрузка & Параметры \\
\hline
Уникальные вставки & \Field{команды, длины, порядок, seed} \\
Поиск & \Field{размер словаря, доля успешных запросов} \\
Удаление & \Field{порядок, число успешных удалений} \\
Смешанные команды & \Field{доли операций, размер словаря, seed} \\
Save и Load & \Field{размер файла и старого дерева, режим кэша} \\
Граничные случаи & \Field{длина 256, крайние числа, ветви удаления} \\
Повреждённые файлы & \Field{повреждение и проверка прежнего словаря} \\
\hline
\end{tabularx}

\Field{Описать валидатор инвариантов, эталон, генерацию допустимых слов
и воспроизведение сбоев. Указать число повторов, медиану и разброс,
одинаковые условия для версий, учёт ввода-вывода и файлового кэша.}
